\chapter{Introduction}

\section{Overview of the Current State of Technology}

In today’s modern world, technology is evolving at a rapid pace all around the world. With this, research and development have produced countless innovations and ideas. Along with the emergence of these inventions and innovations is the concept of ownership of the ideas behind them, referred to as the “creation of the mind” or intellectual property (IP). The “incorporeal right” or the intellectual property right (IPR) is independent from the concept of ownership of the actual object; where owning the tangible object does not guarantee ownership of its IP (Bentley et al., 2020). International organizations such as the World Intellectual Property Organization (WIPO) function with the concept of IPR ownership and pioneer the promotion and protection of IPs across borders all over the world.

The Philippines recognizes the importance of science and technology and their role in research and development, invention, and innovation with the implementation of the Philippine Technology Transfer Act of 2009 (RA 10055). This act also gives emphasis to the generation, utilization, and transfer of intellectual property. It mandates Research and Development Institutions (RDI) to establish Technology Licensing Offices (TLOs) and/or Technology Business Development Offices wherein these institutions have their own policies for technology transfer and IP management. One of these institutions is the Technology Transfer and Business Development Office (TTBDO) of the University of the Philippines Visayas (UPV).

The TTBDO of UPV is responsible for overseeing, managing, protecting, and commercializing the entire UPV IP portfolio, which includes patents, utility models, copyrights, industrial designs, and trademarks. The current IP management system has preliminary storage of IPs on the local machines of the staff. The staff is also responsible for noting important dates and deadlines, such as the timely payment of securing rights and application expiry. Important documents and protected IPs are stored in third-party cloud storage services such as Google Drive. The institution tracks the IP portfolio using online spreadsheet software such as Google Sheets. Through the Google Sheets monitoring mechanism, the TTBDO derives reports through built-in cell functions, which show a tabulated summary of the portfolio.

In light of the growing volume and complexity of IP activities at UPV, the present study implements IRIS (Intellectual Rights Information System), a UPV IP Management Information System. IRIS is a web-based management information system that is designed to streamline the information management process of the Technology Transfer and Business Development Office of the University of the Philippines Visayas. It is also designed to be scalable for the usage of other university campuses that are constituents of the University of the Philippines System.

\section{Problem Statement}

With extensive collaboration and research on the internal processes of the TTBDO, several issues and inefficiencies arising from the current system were observed. The institution performs tracking of the assets in the portfolio by utilizing Google Sheets and storing files in local machines. This poses challenges such as high dependency on the personal local machine, volatility of information upon editing spreadsheet cells, and difficulty in tracking current IP statuses. Because of the nature of the current management system, the staff are extremely cautious when it comes to handling the portfolio; even going to the lengths of downloading a local copy of the recent version of the repository for backup. Additionally, the automated report system is not robust, as it lacks a proper dashboard for analytics and only utilizes the built-in functions of the spreadsheet software, such as “count-if”. The lack of a centralized database increases the risks of data loss, duplication, and miscommunication among staff and stakeholders. Without a structured system, it may also be difficult to monitor and generate summarized descriptive reports that are crucial for decision-making for university officials of UPV and other UP Constituent Universities. Since IPOPHL imposes a maximum two-month response period for ongoing applications, the absence of a reliable notification or alert system increases the risk of application lapses or case withdrawals, which wastes institutional resources and may discourage technology generators from pursuing protection for their research outputs. Furthermore, there is no public portal for showcasing published intellectual property applications, limiting visibility to potential partners, stakeholders, and the general public.

\subsubsection{Summary of Challenges}

The existing processes of the UPV TTBDO mainly require meticulous manual labor, which leads to high risks of errors and mistakes; not only does this greatly affect the integrity of the UPV TTBDO, but it also risks massive damages to intellectual property rights (IPR). These challenges highlight the need for an IP management information system with the following features:

\begin{itemize}
  \item A secure and centralized IP information record storage that is accessible to authorized users.
  \item A robust mechanism for managing IP information.
  \item A monitoring system for up-to-date visibility of multiple IP statuses.
  \item A capacity to generate a comprehensive summary, chart, and descriptive statistics of the IP portfolio for monitoring, compliance, and decision-making.
  \item A user management system to ensure data confidentiality, proper assignment of accountability, and role-based access control (RBAC) with users such as TTBDO Admins, Technology Generators, UP Officials, and Guests.
  \item A system where Technology Generators can submit applications, upload application-related documents, and monitor their application within the system.
  \item A public portal for published IP applications to enhance transparency and facilitate partnerships with external stakeholders.
  \item An automatic alert or notification system for upcoming deadlines or required actions such as IPOPHL’s two-month application period.
\end{itemize}

While digital information management systems exist in global and national agencies such as the World Intellectual Property Organization (WIPO) and the Intellectual Property Office of the Philippines (IPOPHL), there is a lack of research and implementation of IP management systems that are specifically designed for academic or university-level institutions in the Philippines. University-level technology transfer offices typically have small teams with overlapping responsibilities, often with limited budgets for acquiring and maintaining enterprise-grade software for their operations. Moreover, it is challenging to acquire software without incurring intensive customization, financial costs, or process reengineering. There are also frameworks and solutions for these problems that exist in global and national agencies, which can serve as references, but would still require adaptation to fit the operational technicalities of Technology Transfer Offices (TLO), such as the TTBDO of the UPV. These challenges urged the need for the development of a digital information management system to streamline the IP management processes of the TTBDO.

\section{Research Objectives}

\subsection{General Objective}

This study aims to design, develop, and evaluate IRIS, a web-based system that eliminates manual IP application tracking inefficiencies, enhances transparency through status monitoring, offers a standardized platform for Technology Generators regarding IPR application-related processes, and provides comprehensive analytics for improved decision-making in IP management and output in the university.

% The research output is a fully functional, secure, and scalable IRIS with role-based access control, application status tracking, automated notification systems, dashboards, and public portals for published IP applications. Its system architecture is designed to be scalable for potential adaptation by other UP constituent universities. It was evaluated through usability testing, penetration testing, and end-to-end testing.

\subsection{Specific Objectives}

To accomplish the general objective of developing the system within the academic year 2025--2026, the specific objectives are defined as follows:

\begin{enumerate}
  \item analyze existing IP management systems to identify best practices and design patterns and develop the system architecture and database schema to address inefficiencies in TTBDO’s manual process, 
  \item develop a centralized system with secure user authentication, role-based access control, application status tracking, and digital submission workflows for guided IP-type selection, disclosure form submission, and document and attachment management,
  \item implement monitoring, communication, and transparency features, including dashboards, automated notifications, and a public portal for published IP applications, and
  \item evaluate the system through end-to-end testing, performance testing, penetration testing, and user acceptance testing, and document the development process, design, and evaluation results to support replicability for other UP constituent universities.
\end{enumerate}

\section{Scope and Limitations of the Research}

The system is developed and evaluated within the academic year 2025--2026. It encompasses four user roles, namely TTBDO Admin, Technology Generator, UP Official, and Guest. Core features include user management, timestamped application tracking, dashboard analytics, automated notifications, digital submissions, guided IP-type selection, disclosure form and application-related document submission, document and attachment management, and a public portal for published IP applications.

The system is limited to the IP management needs and context of UPV. While the system architecture is designed to be scalable and adaptable for other UP constituent universities, the development, testing, and evaluation were conducted specifically within the UPV environment. Adaptation and deployment to other institutions may require adjustments in the configuration and remain as potential future work.

The system does not include direct integration with IPOPHL's internal filing systems, as applications must still be manually submitted to the office. Additionally, payment processing integration and automated prior art search functionalities are excluded from the system.

Finally, the post-grant operations of TTBDO, such as licensing, commercialization, and IP portfolio management beyond the granted IPs, are outside the scope of this study. The system focuses on streamlining the application process and monitoring of IP assets up to the point of grant approval.
\section{Significance of the Research}

The development of IRIS holds significant value across the following stakeholders:

\subsubsection*{Academic Community}

This research contributes to the limited body of literature on IP management systems specifically designed for Philippine academic institutions by providing a reference architecture for MIS where full system integration with external agencies such as IPOPHL may not be feasible. The system design and implementation approach can serve as a model for other universities, particularly UP constituent universities, in developing their respective IP management systems or migrating from manual to automated IP management processes.

\subsubsection*{TTBDO}

The system addresses inefficiencies in the IP-related processes of TTBDO by supporting application status tracking and monitoring and centralizing document and attachment management. With these improvements, TTBDO staff can then redirect focus from repetitive administrative tasks to higher-value activities such as investor consultation, all while minimizing risks of errors, omissions, and lost documentation.

\subsubsection*{Technology Generators}

The system streamlines the IP application process by allowing submission and monitoring of applications within the system. Technology Generators can use the guided workflow to identify the appropriate IP application type, submit disclosure forms and application-related documents, and get notified of any required actions regarding their application. This effectively reduces friction in protecting Technology Generators’ inventions by allowing them to process these activities with less dependence on office visits or manual email updates.

\subsubsection*{University Officials}

The analytics and reporting capabilities provide university officials with comprehensive data visualization and statistical reports on the IP portfolio and application pipeline. These tools enable data-driven planning for research priorities, technology transfer investments, and resource allocation. The system's monitoring features also support compliance reporting for the UP system.

\subsubsection*{External Stakeholders and General Public}

The public portal enhances transparency, visibility, and collaboration by providing access to published intellectual property applications. This visibility may help support future partnerships and broader engagement with external stakeholders. By showcasing UPV's intellectual property portfolio, the system also helps strengthen the university's presence within industry and innovation ecosystems.